\section{Proof of the candidate}
\label{sec:p147-candidate-proof}

The new point is a function-level estimate for a fixed birational linear
system.  It uses the actual sup norms, not just the dimension or
multiplier ideals of the section spaces.

If \(\dim X=0\), connectedness makes \(X\) a point.  Every nonzero
section space has normalized evaluation supremum \(0\), so every strict
interior slice and its two envelopes equal \(0\).  The terminal assertion
is immediate as well.  Henceforth assume \(\dim X\geq1\).

\subsection{The fixed-degree estimate}

Let \(W=\bigoplus_{a\geq0}W_a\subseteq R(X,L)\) be a multiplicative graded
linear series and set
\[
 u_a=\sup_{\substack{s\in W_a\\\|s\|_{h^a,\infty}\leq1}}
                   a^{-1}\log|s|_{h^a}^2,\qquad
 u_W=(\sup_a u_a)^*.
\]
Assume \(u_W\not\equiv-\infty\).  Each \(u_a\leq0\), hence \(u_W\leq0\).
Fix a positive integer \(m\) for which the rational map defined by \(W_m\)
is birational onto its image.  We will prove
\begin{equation}\label{eq:p147-fixed-degree-bound}
                         P_\theta[u_m]\leq u_W.
\end{equation}

\paragraph{Step 1: resolve the fixed linear system and its obstacle.}
Take a projective resolution of the base ideal of \(W_m\),
\(\mu:Y\to X\), with \(Y\) smooth and connected.  There are an effective
Cartier divisor \(E\), a globally generated line bundle \(A\), and an
isomorphism
\begin{equation}\label{eq:p147-resolved-bundle}
             \mu^*L^{\otimes m}=A\otimes\mathcal O_Y(E).
\end{equation}
Dividing the pulled-back sections by the canonical section \(s_E\)
gives a basepoint-free subspace \(V\subseteq H^0(Y,A)\).
Let \(f:Y\to Z\subseteq\mathbf P(W_m^*)\) be its morphism.
Then \(f\) is birational, \(Z\) is integral, and
\(A=f^*\mathcal O_Z(1)\).  In particular \(A\) is big as well as globally
generated.  We do not assume that \(Z\) is normal or that \(V\) is the
complete linear system.

Give \(A\) a smooth Fubini--Study metric \(a\) from \(V\);
\(\alpha=c_1(A,a)\) is semipositive.
Use \eqref{eq:p147-resolved-bundle} to define the smooth metric \(h_E\) on
\(\mathcal O_Y(E)\) by \(\mu^*h^m=a\otimes h_E\), and put
\[
 g=\log|s_E|_{h_E}^2,\qquad \chi=-g.
\]
Thus \(\chi\) is a lower semicontinuous function, smooth off \(E\),
bounded below, and equal to \(+\infty\) on \(E\).
Moreover
\[
 \alpha+c_1(\mathcal O_Y(E),h_E)=m\mu^*\theta,\qquad
 \mathrm d\mathrm d^{\mathrm c}g=[E]-c_1(\mathcal O_Y(E),h_E).
\]
There is a continuous bounded function \(v_m\in\operatorname{PSH}(Y,\alpha)\)
such that
\begin{equation}\label{eq:p147-fixed-singularity}
                       m\mu^*u_m=g+v_m.
\end{equation}
Indeed, \(v_m\) is the logarithm of the dual norm of evaluation on the
finite-dimensional normed space \(V\).  Evaluation never vanishes on
\(Y\); compactness gives both bounds.  Equivalence with a Euclidean norm
also shows directly the stated analytic singularity type.

Define
\begin{equation}\label{eq:p147-residual-envelope}
 Q=\left(\sup\{q\in\operatorname{PSH}(Y,\alpha):q\leq\chi\}\right)^*.
\end{equation}
This is a bounded \(\alpha\)-psh function.  For the lower bound, use a
constant below \(\inf_Y\chi\), since \(\alpha\geq0\).
For the upper bound, choose a nonempty relatively compact coordinate ball
outside \(E\), on which \(\chi\) is bounded above.
The usual \(L^1\)-compactness of the normalized \(\alpha\)-psh functions
implies a uniform global upper bound for all \(q\leq\chi\): otherwise,
subtracting unbounded maxima would give normalized functions tending to
\(-\infty\) on that ball, contradicting this compactness.
Consequently the supremum is locally uniformly bounded above.
Its regularization satisfies \(Q\leq\chi\) off \(E\), where \(\chi\) is
continuous; on \(E\) the inequality is automatic.  Thus \(Q\) itself is
an admissible competitor.

The definitions give the following exact identity of psh potentials:
\begin{equation}\label{eq:p147-envelope-identification}
                   \mu^*P_\theta[u_m]=\frac{g+Q}{m}.
\end{equation}
Here are the details, including descent.
If \(v\leq0\) is \(\theta\)-psh and \(v\preceq u_m\), then
\(m\mu^*v-g\) is \(\alpha\)-psh away from \(E\).
By \eqref{eq:p147-fixed-singularity} it is locally bounded above near
\(E\), so its local psh weights extend across \(E\).
The extension is at most \(\chi\), so is at most \(Q\).
Conversely, for any \(\alpha\)-psh \(q\leq\chi\), the potential
\((g+q)/m\) is \(\mu^*\theta\)-psh and is at most \(0\).
It descends to a \(\theta\)-psh function on \(X\).
To see the last assertion without a smoothness assumption on an image
linear system, work on the open set where \(\mu\) is an isomorphism.
The descended local psh weights are bounded above near the exceptional
image, because their pullbacks are bounded above over compact
neighbourhoods.  Removable singularities extend them across that analytic
set.  Pulling back these extensions recovers the original weights:
both are psh and agree outside a proper analytic subset.
Finally \(q\) is bounded above and \(v_m\) is bounded below, so the
descended function is \(\preceq u_m\).
The two constructions identify the competitors.
Their upper envelopes give \eqref{eq:p147-envelope-identification}
first away from the exceptional and divisor sets, then everywhere by
equality of psh functions almost everywhere.

\paragraph{Step 2: quantize this extended-real obstacle.}
For \(k\geq1\), define the genuine norm
\[
 \|s\|_{\chi,k}
   =\sup_Y |s|_{a^k}\exp(kg/2),\qquad s\in H^0(Y,A^{\otimes k}),
\]
and the finite-degree potential
\[
 b_{\chi,k}
  =\sup_{\|s\|_{\chi,k}\leq1}k^{-1}\log|s|_{a^k}^2.
\]
The norm is finite since \(\exp(g)\) is smooth and bounded.
It is definite since it is positive on a nonempty open set and a
holomorphic section vanishing there vanishes identically.
Every \(b_{\chi,k}\) is at most \(Q\).

We claim that
\begin{equation}\label{eq:p147-extended-obstacle-quantization}
                         b_{\chi,k}\longrightarrow Q
                  \quad\hbox{almost everywhere on }Y.
\end{equation}
This only needs the smooth-obstacle Bergman theorem for a big line bundle.
For clarity, set
\[
             \chi_j=-\log\bigl(|s_E|_{h_E}^2+j^{-1}\bigr).
\]
These smooth functions increase to \(\chi\).
Let \(Q_j=P_\alpha(\chi_j)\).
They increase to \(Q\) after upper semicontinuous regularization.
Indeed, for any \(q\leq\chi\) and \(\varepsilon>0\), the increasing open
sets
\(\{q-\varepsilon<\chi_j\}\) cover the compact space \(Y\).
One of them is all of \(Y\), so \(q-\varepsilon\leq Q_j\).
Taking suprema and then regularizations proves the assertion.
In particular \(Q_j\to Q\) almost everywhere.

For each fixed \(j\), use the smooth metric \(a\exp(-\chi_j)\) on the big
line bundle \(A\).
The Bergman-weight convergence in \cite[Theorem~1.4]{Berman07} gives
the full-section \(L^2\) quantization of \(Q_j\), uniformly on compact
subsets of \(Y\setminus\mathbf B_+(A)\).
The auxiliary fixed bundle in that theorem is the trivial bundle.
There is no positivity assumption on this smooth metric.
The same assertion holds for the sup norm: with a fixed smooth volume
form of total mass one,
\[
 \|s\|_{2,j}\leq\|s\|_{\infty,j}
       \leq C_j k^n\|s\|_{2,j}
\]
for all \(k\) sufficiently large.  For the upper estimate cover \(Y\)
by coordinate charts with smooth local weights, apply the holomorphic
mean-value inequality on balls of radius \(k^{-1}\), and note that
\(k\) times the oscillation of a smooth weight on such a ball is bounded.
This gives a polynomial loss, sufficient here.
Consequently the two logarithmic evaluation suprema differ by at most
\(2\log(C_jk^n)/k\), which tends to zero.
Thus, if \(b_{\chi_j,k}\) uses the sup norm associated with \(\chi_j\),
\[
                  b_{\chi_j,k}\longrightarrow Q_j
             \quad\hbox{on }Y\setminus\mathbf B_+(A).
\]
Since \(\chi_j\leq\chi\), their unit balls satisfy
\(b_{\chi_j,k}\leq b_{\chi,k}\leq Q\).
Taking a lower limit in \(k\), then the supremum over the countable set
of \(j\), proves \eqref{eq:p147-extended-obstacle-quantization}.
The exceptional set is the augmented base locus and the null set where
the increasing \(Q_j\)'s differ from their regularized limit.
Neither uniform convergence at \(E\) nor continuity of \(Q\) is used.

\paragraph{Step 3: return complete sections to the given graded series.}
This is where birationality matters.
Let \(\nu:Z^\nu\to Z\) be the normalization.  It is finite.
As \(Y\) is normal and \(f\) is proper birational,
\[
                         f_*\mathcal O_Y=\nu_*\mathcal O_{Z^\nu}.
\]
For example, factor \(f\) through \(Z^\nu\).
A function regular upstairs is a rational function on the normal target
without a pole at any prime divisor, hence is regular; this gives the
displayed identity.  Alternatively apply \cite[Tag~0AY8]{StacksBirational}:
the morphism \(Y\to Z^\nu\) is proper, the target is normal and integral,
the source is integral and reduced, and its generic fibre is the single
point with function field \(\mathbf C(Z^\nu)\).

The coherent quotient
\(\nu_*\mathcal O_{Z^\nu}/\mathcal O_Z\) has zero generic stalk.
Its annihilator \(\mathfrak c\) is therefore a nonzero coherent ideal and
\[
              \mathfrak c\,\nu_*\mathcal O_{Z^\nu}\subseteq\mathcal O_Z.
\]
Since \(\mathcal O_Z(1)\) is ample, choose a fixed integer \(r\geq1\) and
a nonzero section
\[
                  b\in H^0(Z,\mathfrak c\otimes\mathcal O_Z(r)).
\]
Global generation of a sufficiently high twist guarantees its existence.
The projection formula gives, for every \(k\),
\begin{equation}\label{eq:p147-conductor-inclusion}
 b\,H^0(Y,A^{\otimes k})
                    \subseteq H^0(Z,\mathcal O_Z(k+r)),
\end{equation}
where pullback to \(Y\) is understood on the left.
For \(k\) large, Serre vanishing applied to the ideal of
\(Z\subseteq\mathbf P(W_m^*)\) makes the restriction of homogeneous
polynomials of degree \(k+r\) onto the right side surjective.
Therefore, for every \(s\in H^0(Y,A^{\otimes k})\),
\begin{equation}\label{eq:p147-polynomial-return}
                 s_E^{\,k+r}bs
          \ \in\ \mu^*\operatorname{im}
                    \bigl(\operatorname{Sym}^{k+r}W_m
                                      \longrightarrow W_{m(k+r)}\bigr).
\end{equation}
This inclusion has the indicated direction.  An arbitrary section of the
normalization does not descend without the conductor factor.
The standard algebraic inputs here are finite normalization and Serre
vanishing/global generation; see \cite{StacksNormalization,StacksSerre}.
No finite generation of the entire \(W\) is involved.

Regard \(b\) now as its pullback, a nonzero section of \(A^{\otimes r}\), and
set
\[
 B=\sup_Y |b|_{a^r}\exp(rg/2),\qquad
 \beta=\log|b|_{a^r}^2.
\]
The constant satisfies \(0<B<\infty\).
If \(\|s\|_{\chi,k}\leq1\), the section in
\eqref{eq:p147-polynomial-return} has \(\mu^*h^{m(k+r)}\)-sup norm at most
\(B\), since its pointwise norm is exactly
\[
       \bigl(|b|_{a^r}\exp(rg/2)\bigr)
       \bigl(|s|_{a^k}\exp(kg/2)\bigr).
\]
After division by \(B\) it is an allowed section in the definition of
\(u_W\).  Taking the supremum over \(s\) yields, away from \(E\cup\{b=0\}\),
\begin{equation}\label{eq:p147-pointwise-return}
 \mu^*u_W\geq
       \frac{g}{m}
       +\frac{k b_{\chi,k}+\beta-2\log B}{m(k+r)} .
\end{equation}
All degrees, norms, and filtration memberships in this inequality are
exact.  Only its limit discards the fixed cost \(r\) and the fixed
logarithmic term.
By \eqref{eq:p147-extended-obstacle-quantization}, its right side converges
almost everywhere to \((g+Q)/m\).
Using \eqref{eq:p147-envelope-identification} gives
\(u_W\geq P_\theta[u_m]\) almost everywhere on \(X\).
Both functions are \(\theta\)-psh, so the inequality holds everywhere.
This proves \eqref{eq:p147-fixed-degree-bound}.

\subsection{Apply the estimate to the filtration}

Fix \(\tau<c\) and take \(W_a=W_a^\tau\).
The series contains an ample series; consequently there is an integer
\(m_0\geq1\) such that the maps defined by \(W_{m_0q}^\tau\) are birational
for all sufficiently large \(q\).
This is the precise input stated in the proof of
\cite[Proposition~4.6]{DXZ23}, referring to \cite[Lemma~1.6]{BC11}.
For \(m=m_0 2^j\) sufficiently large,
\eqref{eq:p147-fixed-degree-bound} gives
\[
                    P_\theta[u_{\tau,m_0 2^j}]\leq u_\tau.
\]
Multiplicativity and the sup-norm constraint give
\(u_{\tau,2^j}\leq u_{\tau,m_0 2^j}\): take the \(m_0\)-th power of any
normalized section.  Monotonicity of the singularity envelope therefore
gives
\begin{equation}\label{eq:p147-dyadic-bound}
                         P_\theta[u_{\tau,2^j}]\leq u_\tau.
\end{equation}
The proof of \cite[Lemma~4.5]{DXZ23} supplies the exact identity
\[
 P_\theta[u_\tau]
      =\left(\lim_{j\to\infty}P_\theta[u_{\tau,2^j}]\right)^*
      =P_\theta[u_\tau]_{\mathcal I}.
\]
In that proof the dyadic potentials increase almost everywhere to
\(u_\tau\); the maximized finite-level potentials are \(\mathcal I\)-model,
and positive mass makes the increasing-limit closure theorem applicable.
It is not an assertion that the raw potentials were already maximal.
Taking the upper semicontinuous regularization in
\eqref{eq:p147-dyadic-bound} gives \(P_\theta[u_\tau]\leq u_\tau\).
The reverse inequality follows from \(u_\tau\leq0\), so
\eqref{eq:p147-target} follows.

For completeness, positive mass on this whole strict interior is not an
extra restriction.  Choose \(a<\tau<b<c\) with \(u_a=V_\theta\), using
linear boundedness.  Write \(\tau=(1-t)a+tb\), where \(0<t<1\).
Concavity gives
\(u_\tau\geq(1-t)V_\theta+t u_b\).
Positivity and multilinearity of nonpluripolar products, together with
mass monotonicity under decreasing singularities, give
\[
       \int_X\theta_{u_\tau}^{\,n}
                    \geq(1-t)^n\int_X\theta_{V_\theta}^{\,n}>0.
\]
This also explains the positive-mass observation immediately before
DXZ Lemma~4.5.  The candidate theorem is proved.

\subsection{The particular terminal convention}

The following observation is separated from the theorem's interior target.
DXZ defines \(u_c=\lim_{\tau\uparrow c}u_\tau\).
If this limit is not identically \(-\infty\), it is a decreasing limit of
psh functions and hence is psh.
For every \(\tau<c\), monotonicity gives
\[
 P_\theta[u_c]\leq P_\theta[u_\tau]=u_\tau,\qquad
 P_\theta[u_c]_{\mathcal I}\leq
                          P_\theta[u_\tau]_{\mathcal I}=u_\tau.
\]
Letting \(\tau\uparrow c\), and using the reverse inequalities from
\(u_c\leq0\), shows that this particular endpoint is also fixed by both
envelopes, even if its top mass is zero.
If \(u_c\equiv-\infty\), it is simply the extended-real terminal slice.
For \(\tau>c\) the slice is defined to be \(-\infty\).
These observations contain no reconstruction statement for unspecified
zero-mass filtration data and do not decide the separate inventory atom
P14.8 without first correcting its exact formulation.
